%% ============================================================
\section{Preliminaries}\label{sec:preliminaries}
%% ============================================================

\subsection{Loop Verification}

A verification task consists of a loop
\(\mathcal L=(P,B,S)\) and a postcondition \(Q\), where \(P\) is the
precondition, \(B\) is the loop guard, and \(S\) is the loop body.  Let
\(\operatorname{Reach}(\mathcal L)\) denote the states reachable at the loop
head, including the final guard test at loop exit.

An invariant \(I\) is \emph{sound} (equivalently, inductive) when it is
established by \(P\) and preserved by one iteration:
\[
  P\Rightarrow I,
  \qquad
  \{I\wedge B\}\,S\,\{I\}.
\]
These obligations ensure that \(I\) holds on every state in
\(\operatorname{Reach}(\mathcal L)\)
~\citep{floyd1967assigning,hoare1969axiomatic}.

We say that \(I_1\) is \emph{stronger} than \(I_2\), written
\(I_1\succeq I_2\), when
\[
  \models I_1\Rightarrow I_2.
\]
An invariant \(I^\star\) is \emph{exact} when
\[
  \denote{I^\star}=\operatorname{Reach}(\mathcal L).
\]
Because every sound invariant contains all reachable states, an exact
invariant is stronger than every sound invariant of the same loop:
\[
  I^\star\text{ exact},\ J\text{ sound}
  \quad\Longrightarrow\quad
  \models I^\star\Rightarrow J.
\]
Exactness is a semantic ideal; the reachable states need not have a finite
representation in the annotation fragment used by the generator.

Finally, a sound invariant \(I\) is \emph{\(Q\)-sufficient} when it proves the
postcondition at loop exit:
\[
  I\wedge\neg B\Rightarrow Q.
\]
Thus exactness describes the loop independently of a target, whereas
\(Q\)-sufficiency is relative to one postcondition.  The final verifier must
prove both soundness obligations and this exit obligation.

\subsection{Target-Hidden Verification}

In \emph{target-hidden loop verification}, the generator and all intermediate
feedback see \(\mathcal L=(P,B,S)\) but not \(Q\).  The postcondition is
restored only after the generated invariant has been fixed.  Preconditions and
executable loop semantics remain visible, while assertions, postconditions,
assertion calls, and conventional error-label names are masked.  The untouched
program is used only for final verification.

\subsection{Scope and Motivating Example}

We focus on numerical loops for two reasons.  First, their invariants are
predominantly arithmetic equalities and inequalities, including polynomial
relations.  These predicates are comparatively amenable to automatic checking,
which provides a scalable and objective verification signal.  Second, numerical
loops factor out the representational machinery of arrays, heaps, and recursive
data structures, whose verification often requires auxiliary functions or
inductive predicates.  This gives a controlled setting in which success depends
more directly on discovering algebraic relations and reasoning about loop
dynamics.

Concretely, we study one braced \texttt{while} loop over scalar C integers,
including nonlinear updates; arrays, heaps, and recursive structures are out
of scope.  This restriction does not make invariant discovery trivial: even a
small loop can hide a nonlinear relation across coupled recurrences.

Consider the following abridged program:
\begin{lstlisting}
int x = 1, y = 1, c = 1;
while (c < k) { c++; x = x*z + 1; y = y*z; }
/*@ assert 1+x*z-x-z*y == 0; */
\end{lstlisting}
The visible assertion can be copied verbatim as an invariant; when hidden, its
nonlinear relation must instead be recovered from the coupled updates.
